\scanpage{76}\mspage{75}

Dans ce cas, l'axiome $C6$ (pour une relation \inserted{d'ordre} suppl $\preceq$ sur $\mathcal F=\unsure{\mathrm P|\mathcal P}_{f}(M)$) signifie \unsure{seulement|seulmt} que si $F'\preceq F$, $F'$ est vide ssi $F$ l'est. Prenons p.ex. \deleted{$\unsure{F|K}\preceq$} $\preceq$ \deleted{\unsure{égal|égale}} à la relation d'ordre discrète, c'est OK\unsure{.|,} \deleted{$C7_{\unsure{0|a}}$ est \unsure{tautologique}} $C6,C7$ tautologiques. On voit \unsure{donc|bien} qu'il manque un axiome, pour assurer la validité du corollaire. On devine que ce sera la partie \unsure{--)|a)} du \unsure{Const|Cont} $4$ (p. 35), \unsure{jointe} à Const $7$. Donc pour l'instant \unsure{je} ne \unsure{vais|sais} pas utiliser que le préordre $\ll$ soit un ordre.
\margin{Aussi, cela résulte \unsure{aussitôt|aussi} \unsure{de la} \unsure{modification} de $C4$ ancien (page $47$)}

\deleted{\unsure{Je} songe à la \unsure{propté|propriété} \unsure{suiv\textsuperscript{te}|suiv.|suivante}, qu'il faudrait avoir tot ou tard :}

\deleted{\deleted{\unsure{Soient|Soit} $\unsure{Y,|\forall}X\in M$ \unsure{si|avec|il} $Y\ll X$, et} supposons que \deleted{$Y'\leq Y$, $Y'\in\partial X$, $Y\cap\partial X\unsure{=|-}$} \deleted{$\unsure{\exists|\nexists}\,Y'\leq Y$, \unsure{alors|les}} $Y'\leq Y$, $(Y'\in M)$ tels que \deleted{$\unsure{\phi|\partial|\exists}\,Y'\ll\partial X$}}

\deleted{Soient $F$ figure, $F'$ \unsure{ss-figure|s-figure}, $G\ll F$, $G'$ \inserted{$\leq G$} le \unsure{raffint|raffinement} de $F'$ induit par $G$ cf plus \unsure{bas|pas}). \deleted{\unsure{Supposons|supposons} que} \inserted{Alors} $G'=\unsure{\phi|\varnothing}$. \deleted{Alors $G$ est} \unsure{ssi|ss} $G$ disjoint de $F'$.}

Il faut développer la notion de \uline{\unsure{raffinements|raffinement|raffinants} induit}. J'ai envie de \deleted{\unsure{démontrer|Il en résulte|Il en est}} \inserted{\unsure{prouver|posons}} l'énoncé (quitte à rajouter des axiomes) :
